\cotefile{29}
\scanpage{102}
\mspage{101}
% Upper-left margin: thin grey curved line; origin unassigned.
\margin{\figure{\begin{tikzpicture}\draw[gray,line width=.2pt] (0,0) .. controls (1,-.13) and (2.2,-.19) .. (3.2,-.14);\end{tikzpicture}}}
% Upper-right margin: two separate diffuse brown-grey traces with darker centres; no textual value or maker assigned.
\margin{\figure{\begin{tikzpicture}
\fill[brown!30!gray,opacity=.45] (-.018,.02) .. controls (.03,.01) and (.04,-.08) .. (.027,-.17) .. controls (-.01,-.20) and (-.04,-.09) .. cycle;
\fill[black!65] (-.012,-.145) .. controls (.015,-.157) and (.023,-.172) .. (.004,-.181) .. controls (-.018,-.178) and (-.024,-.159) .. cycle;
\fill[brown!30!gray,opacity=.45] (.385,-.66) .. controls (.42,-.72) and (.43,-.83) .. (.408,-.96) .. controls (.39,-1.07) and (.36,-1.14) .. (.348,-1.18) .. controls (.326,-1.19) and (.341,-1.08) .. (.365,-.96) .. controls (.39,-.83) and (.36,-.72) .. cycle;
\fill[black!65] (.372,-.90) ellipse (.015 and .018);
\fill[black!65] (.371,-1.01) .. controls (.391,-1.035) and (.376,-1.06) .. (.356,-1.063) .. controls (.348,-1.045) and (.36,-1.02) .. cycle;
\fill[black!45] (.350,-1.10) .. controls (.373,-1.085) and (.371,-1.13) .. (.341,-1.16) .. controls (.335,-1.145) and (.341,-1.12) .. cycle;
\end{tikzpicture}}}

Examinons maintenant le cas où l'un des cycles $\underline{a}$ ou $\underline{b}$ (par exemple $\underline{b}$) appartient au noyau d'Albanese $Z_{\handwritten{\ell}}(A)$ ; par linéarité, on se ramène au cas où $\underline{b}$ est de la forme $\underline{c}_u-\underline{c}$, avec $\underline{c}\handwritten{\in}Z^{s}_{\unsure{o|0}}(A)_{\bar{\underline{k}}}$, et $u\handwritten{\in}{}^{s}A_{\bar{\underline{k}}}$. D'après l'invariance du symbole $j_{\handwritten{o}}$ par la translation $\handwritten{\tau_u}$, les deux membres de (3) sont respectivement égaux à $j_{\handwritten{o}}(Y_{\underline{a}},\underline{c})$ et à $j_{\handwritten{o}}(Y^{-}_{\underline{c}},\underline{a})$\inserted{\handwritten{, en posant $Y=X_u-X$.}} Puisqu'on a $X\handwritten{\in}D_a(A)$, on a $Y\handwritten{\sim}0$, et on est ramené au cas déjà traité.
% The handwritten insertion begins with comma + en at the right end of line 6;
% posant Y=X_u-X. continues in the left margin at line 7, before Puisqu'on.

Supposons maintenant $X$, $\underline{a}$ et $\underline{b}$ quelconques. On peut, en modifiant $\underline{b}$ par l'addition d'un élément de $Z^{\handwritten{s}}_{\handwritten{\ell}}(A)_{\handwritten{\bar{\underline{k}}}}$, se ramener au cas où le cycle $\underline{a}^{-}\handwritten{\otimes}\underline{b}$ est étranger à $X$. Les symboles $X_{\underline{a}}(\underline{b})$ et $X^{-}_{\underline{b}}(\underline{a})$ sont alors définis, et on a
\[
X_{\underline{a}}(\underline{b})=X^{-}_{\underline{b}}(\underline{a}),\quad\text{d'où}\quad v(X_{\underline{a}}(\underline{b}))=v(X^{-}_{\underline{b}}(\underline{a})).
\]

D'autre part, soit $a$ (resp. $b$) l'un quelconque des composants de $\underline{a}$ (resp. $\underline{b}$). Alors les symboles $i(X_a,b)$ et $i(X^{-}_b,a)$ sont définis ; l'application $\handwritten{\unsure{\zeta|\xi}}:A\handwritten{\longrightarrow}A$ définie par $\handwritten{\zeta}(x)=a+b-x$ étant \uline{p}-isomorphique en $b^{\handwritten{\unsure{\circ|o|0}}}$, de valeur $a^{\handwritten{\circ}}$, on a, d'après la propriété d'invariance du symbole $i$,
\[
i(X_a,b)=i(X^{-}_b,a).
\]

Il en résulte, par linéarité,
\[
i(X_{\underline{a}},\underline{b})=i(X^{-}_{\underline{b}},\underline{a}).
\]
% Left margin beside the second i display: short dark tapering stroke, origin unassigned.
\margin{\figure{\begin{tikzpicture}\fill (0,0) .. controls (.012,-.06) and (.018,-.14) .. (.019,-.23) .. controls (.006,-.24) and (.002,-.19) .. (.004,-.13) .. controls (-.004,-.07) and (-.008,-.02) .. cycle;\end{tikzpicture}}}

Compte tenu des formules (2) précédentes, on en déduit (3).

Notons enfin que $j_{\handwritten{o}}(X,\underline{a})$ s'annule lorsque tous les composants de $A$ appartiennent au groupe ${}^{s}A_{\bar{\underline{k}},\handwritten{o}}$. En effet, soit $f$ une fonction sur $V$, définie sur $\underline{k}$, telle que

% archive sheet number: 101
